\documentclass[12pt,a4paper]{ctexart}

\usepackage[margin=2.4cm]{geometry}
\usepackage{amsmath}
\usepackage{booktabs}
\usepackage{siunitx}
\usepackage{hyperref}
\usepackage{enumitem}
\usepackage{xcolor}

\hypersetup{colorlinks=true, linkcolor=blue, urlcolor=blue}
\sisetup{separate-uncertainty=true, per-mode=symbol}

\title{\textbf{布朗运动测量玻尔兹曼常数实验\\误差与偏移项分析说明}}
\author{Steven Chen}
\date{2026 年 5 月}

\begin{document}
\maketitle

\section{说明目的}

本文用于向组员说明当前数据处理中发现的主要误差来源，尤其是布朗运动轨迹中可能存在的慢漂移项对均方位移（MSD）和玻尔兹曼常数 $k_B$ 测量的影响。

结论先行：当前报告中的数值计算本身可以复现，但第二、第三、第四颗粒子的原始布朗轨迹存在明显的非理想成分。若直接用全轨迹 MSD 作为参考，粒子 2、3、4 的扩散系数会分别被抬高约 $23\%$、$25\%$ 和 $120\%$；粒子 5 的全轨迹扩散系数仅偏高约 $4.7\%$，与粒子 1 较接近。因此，当前数据更适合描述为“存在漂移污染的数据分析”，不宜简单表述为完全理想的布朗运动测量。

\section{理想模型与偏移项}

理想的一维布朗运动满足
\begin{equation}
  \langle \Delta x \rangle = 0, \qquad
  \langle \Delta x^2(\tau) \rangle = 2D\tau .
\end{equation}

实验中观测到的位置往往不只是纯布朗运动，可写成
\begin{equation}
  x_{\mathrm{obs}}(t)=x_{\mathrm{B}}(t)+x_{\mathrm{drift}}(t)+\epsilon(t),
\end{equation}
其中 $x_{\mathrm{B}}(t)$ 是布朗运动，$x_{\mathrm{drift}}(t)$ 是慢漂移，$\epsilon(t)$ 是定位噪声。

若漂移近似为匀速 $v_{\mathrm{d}}$，则位移为
\begin{equation}
  \Delta x_{\mathrm{obs}} = \Delta x_{\mathrm{B}} + v_{\mathrm{d}}\tau .
\end{equation}
于是均方位移变为
\begin{equation}
  \langle \Delta x_{\mathrm{obs}}^2 \rangle
  \approx 2D\tau + (v_{\mathrm{d}}\tau)^2 .
\end{equation}

这个额外的 $(v_{\mathrm{d}}\tau)^2$ 项会使 MSD 偏大。更一般地，若轨迹中存在气流、电场不稳定、液滴缓慢移动或局部非平稳运动，也会在长时间尺度上抬高 MSD。

\section{当前实验中的显著误差}

当前五颗粒子的全轨迹参考扩散系数 $D_{\mathrm{ref}}$ 与理论 Stokes--Einstein 期望值 $D_{\mathrm{expect}}$ 的比较如下：

\begin{table}[h]
\centering
\caption{全轨迹参考扩散系数与理论值对比}
\begin{tabular}{lccc}
\toprule
粒子 & $D_{\mathrm{ref}}$ ($\si{\micro\metre^2\per\second}$) & $D_{\mathrm{expect}}$ ($\si{\micro\metre^2\per\second}$) & 比值 \\
\midrule
粒子 1 & 19.03 & 18.68 & 1.019 \\
粒子 2 & 20.91 & 16.95 & 1.234 \\
粒子 3 & 19.12 & 15.29 & 1.250 \\
粒子 4 & 35.09 & 15.97 & 2.198 \\
粒子 5 & 19.52 & 18.64 & 1.047 \\
\bottomrule
\end{tabular}
\end{table}

可以看到，粒子 1 与粒子 5 与理论值基本接近；粒子 2、3、4 的 $D_{\mathrm{ref}}$ 明显偏大，分别约为理论值的 $1.23$ 倍、$1.25$ 倍和 $2.20$ 倍。

这意味着：若直接把全轨迹 MSD 当作理想布朗运动参考，粒子 2、3、4 的 $\langle x^2\rangle$ 会系统性偏大。若用这些偏大的 MSD 计算 $k_B$，也会得到明显偏大的 $k_B$。粒子 5 的全轨迹参考值则相对更稳定。

\section{窗口选择带来的方法问题}

当前程序中使用了目标优化窗口选择策略，即通过选择局部区间，使计算出的 $k_B$ 接近预设目标。这个做法可以找到“数值上较好”的局部窗口，但它不是完全盲法测量。

独立复算显示：若固定当前被选中的窗口，$k_B$ 的数值计算是正确的，误差约为 $+3\%$。但是，如果只按统计质量评分选窗口，受漂移影响较大的粒子会得到明显偏高的 $k_B$ 误差。

因此当前结果应谨慎表述为：在轨迹中选取了局部较合适的时间区间后，得到接近理论值的结果。不能简单表述为整条原始轨迹都与理想布朗运动高度一致。

\section{可能导致误差的原因}

根据数据表现，主要可能原因包括：

\begin{enumerate}[leftmargin=2em]
  \item \textbf{悬浮不稳定或气流扰动：} 油滴在悬浮时可能受到微弱空气流动或热对流影响，产生低频定向运动。
  \item \textbf{电场或仪器漂移：} 密立根装置中的电场不均匀、平台振动或成像系统漂移可能导致油滴位置缓慢移动。
  \item \textbf{有限轨迹长度与非平稳性：} 粒子 2 和粒子 3 的布朗轨迹长度明显短于粒子 1，统计平均不充分；粒子 4 虽然轨迹较长，但全局漂移更强；粒子 5 轨迹较长且全轨迹扩散系数更接近理论值。
  \item \textbf{单窗口样本数不足：} 当前每个窗口使用 $N=150$ 个独立位移，理论相对标准误差约为
  \begin{equation}
    \sqrt{\frac{2}{N}} = \sqrt{\frac{2}{150}} \approx 11.5\% .
  \end{equation}
  因此单个窗口天然会有 $10\%$ 量级的随机波动。
  \item \textbf{长时间 MSD 对漂移敏感：} 用较长 lag 拟合 MSD 时，漂移项 $(v_{\mathrm{d}}\tau)^2$ 的影响会增强。
\end{enumerate}

\section{关于慢漂移扣除的说明}

一种尝试是对轨迹做低频去趋势，例如用 $80\,\mathrm{s}$ 滚动平均估计慢漂移：
\begin{equation}
  x_{\mathrm{slow}}(t)=\mathrm{rolling\ mean}_{80\mathrm{s}}[x(t)],
\end{equation}
再构造
\begin{equation}
  x_{\mathrm{corr}}(t)=x(t)-x_{\mathrm{slow}}(t)+\overline{x} .
\end{equation}

这个方法可以明显降低粒子 2、3、4 的 MSD 偏差，使 $D_{\mathrm{ref}}$ 更接近理论值。但它有重要限制：$x_{\mathrm{slow}}(t)$ 是从油滴自身轨迹中估计出来的，可能同时扣除了真实布朗运动中的低频涨落。因此它更适合作为“漂移敏感性分析”，不应直接作为严格的主实验结论。

更中立的表述是：低频去趋势说明当前结果对慢漂移较敏感；去趋势后数值改善，支持“原始轨迹中存在低频漂移污染”的判断。

\section{建议的解决方案}

\subsection{数据分析层面}

\begin{enumerate}[leftmargin=2em]
  \item \textbf{报告原始轨迹结果与去趋势结果两套结果：} 原始结果代表直接测量；去趋势结果作为漂移校正后的参考。
  \item \textbf{缩短参考 MSD 拟合范围：} 例如比较 $\tau \le 2\,\mathrm{s}$、$5\,\mathrm{s}$、$10\,\mathrm{s}$、$20\,\mathrm{s}$ 的 $D_{\mathrm{ref}}$，避免长 lag 被漂移放大。
  \item \textbf{做分段稳定性检查：} 每隔固定时间计算局部 $D$、局部平均位移和漂移斜率，判断哪些时间段不稳定。
  \item \textbf{避免用 NIST 值直接选窗作为主结论：} 若使用目标优化窗口，应明确说明这是后验选择或校准选窗。
  \item \textbf{给出统计不确定度：} 对 $N=150$ 的单窗口结果，应说明理论随机误差约 $11.5\%$，不能声称单窗口自然达到 $2\%$--$3\%$ 精度。
\end{enumerate}

\subsection{实验采集层面}

\begin{enumerate}[leftmargin=2em]
  \item \textbf{延长采集时间：} 若希望单窗口统计误差达到 $3\%$，理论上需要约 $N\approx 2222$ 个独立位移，远大于当前 $N=150$。
  \item \textbf{等待悬浮稳定后再记录：} 减少开始阶段的电场调节和油滴松弛过程。
  \item \textbf{控制气流和温漂：} 尽量封闭观察区域，减少热对流和外界扰动。
  \item \textbf{记录参照物或多颗粒子：} 若能同时记录固定背景点或多颗粒子，可用公共漂移进行更物理可靠的漂移扣除。
  \item \textbf{记录二维坐标：} 同时使用 $x$ 和 $y$ 的布朗涨落可提高统计量，但需避免竖直方向受电场或重力影响。
\end{enumerate}

\section{总结}

当前实验的主要问题不是公式计算错误，而是粒子 2、3、4 的原始轨迹中存在明显低频非理想运动，使全轨迹 MSD 偏大。粒子 5 的全轨迹结果较接近理论值，可作为相对稳定的一组补充数据。慢漂移扣除可以改善数值，但由于它从油滴自身轨迹中估计漂移，可能同时滤除真实布朗运动的低频成分，因此只能作为辅助分析。

对组内讨论，建议采用如下表述：

\begin{quote}
当前数据表明粒子 2、3、4 的布朗轨迹受到慢漂移或非平稳运动影响。原始全轨迹 MSD 给出的扩散系数分别偏高约 $23\%$、$25\%$ 和 $120\%$；粒子 5 偏高约 $4.7\%$，与理论更接近。通过局部选窗或低频去趋势可以得到更接近理论的结果，但这类处理需要明确说明其后验性和模型假设。若要获得更可靠的独立测量，应改进采集稳定性、延长轨迹长度，并使用独立参照物或多粒子数据估计漂移。
\end{quote}

\end{document}
